% SOURCE_STATEMENT: "For coupled magnetomechanical systems, should relaxation be formulated monolithically or as a staggered sequence of partial relaxations, and under what conditions are the two procedures equivalent?"
% Preparation and final self-audit: CONTRIBUTING.md prompts C and D; source/audit artifacts retained in the local physics_statements_2026_10_11 batch.
% !TeX program = lualatex
\documentclass[11pt]{article}
\usepackage[margin=25mm]{geometry}
\usepackage{fontspec}
\setmainfont{Latin Modern Roman}
\usepackage{amsmath,amssymb,mathtools}
\usepackage{unicode-math}
\setmathfont{Latin Modern Math}
\usepackage{xurl,hyperref}
\hypersetup{colorlinks=true,urlcolor=blue}
\setcounter{secnumdepth}{0}
\setlength{\parindent}{0pt}
\setlength{\parskip}{5pt}
\setlength{\emergencystretch}{3em}
\newcommand{\sref}[2]{\hyperlink{source:#1:#2}{[S#2]}}
\newcommand{\cataloguescope}[1]{\paragraph{Recorded target.}#1}
\newcommand{\R}{\mathbb R}
\newcommand{\C}{\mathbb C}
\newcommand{\N}{\mathbb N}
\newcommand{\Z}{\mathbb Z}
\newcommand{\Q}{\mathbb Q}
\newcommand{\F}{\mathbb F}
\newcommand{\eps}{\varepsilon}
\begin{document}
% UnsolvedMath ID 30004442; source catalogue code OWR-17476-002
\setcounter{section}{30004441}
\section{Monolithic versus Staggered Magnetomechanical Relaxation}\label{30004442}
{\small UnsolvedMath ID 30004442 · Mathematical Physics}
\subsection{Definitions and mathematical statement}

\paragraph{Shared notation.}$\mathbb N=\{1,2,\ldots\}$; the standard number fields have their usual meanings. All additional conventions are stated locally.


\paragraph{Definitions and conventions.}
Relaxation replaces a nonconvex material energy by an effective energy allowing fine-scale microstructure, commonly through a lower-semicontinuous envelope or an admissible microstructure minimization. Monolithic relaxation treats mechanical and magnetic variables together. Staggered relaxation alternates or successively relaxes subsets of those variables. Equivalence requires equality of the resulting effective energies or solution behavior under stated coupling and admissibility assumptions.

\paragraph{Target retained from the canonical source catalogue.}
\paragraph{Catalogue wording retained for traceability.}
For coupled magnetomechanical systems, should relaxation be formulated monolithically or as a staggered sequence of partial relaxations, and under what conditions are the two procedures equivalent?

\paragraph{Source scope and qualification.}
The ordering, topology and compatibility constraints of the partial relaxations matter. Equality of two numerical minimization runs does not establish equivalence of the variational relaxation procedures.
\subsection{Short English statement}
Decide when combined and successive relaxations give the same effective magnetomechanical model.
\subsection{Sources}
\begin{itemize}
\item[\hypertarget{source:30004442:1}{S1}] ULAM.AI and the UnsolvedMath Contributors, UnsolvedMath, record 30004442 (OWR-17476-002). Catalogue curation under CC BY 4.0; definitions and English exposition added locally.
\url{https://huggingface.co/datasets/ulamai/UnsolvedMath}.
\item[\hypertarget{source:30004442:2}{S2}] Mechanics of Materials: Towards Predictive Methods for Kinetics in Plasticity, Fracture, and Damage (2020), Oberwolfach Reports, DOI 10.4171/owr/2020/13, pp. 713--783. Primary report accessed 2026-10-11; relevant PDF pages 27.
\url{https://doi.org/10.4171/owr/2020/13}.
\item[\hypertarget{source:30004442:3}{S3}] Rossi--Thomas (2017) (catalogue-listed context; not independently consulted in this preparation).
\url{https://doi.org/10.1137/15M1051567}.
\item[\hypertarget{source:30004442:4}{S4}] Crismale--Lazzaroni--Rossi (2023) (catalogue-listed context; not independently consulted in this preparation).
\url{https://doi.org/10.1007/s10231-022-01280-0}.
\end{itemize}
\label{end:30004442}
\end{document}
